\section{部署与观测}

\subsection{监控 inference-time 指标}
Chen 把 inference-time pipeline 看成一个“黑盒”——输入 prompt/block 之后，依次走过 sampling、search、trace，每一段都应该输出 metrics。典型的 metrics 包含 sampling depth、branching ratio、verifier score distribution、trace iteration count 与 latency。建议在 dashboard 上按 5 分钟一组绘制这些曲线，真实系统中一旦 sampling depth 突然下降或 trace iteration 突增，就要触发调查。

\begin{knowledgebox}{Observability 工具链}
1. Sampling log：记录 prompt + generated candidates + token count；2. Search log：记录 beam/tree status (depth, score, branch id)；3. Trace log：输出每轮 self-correction 是否有人类命名的 error pattern；4. Alert rules：设定 collector，在 sample 多样性低于 baseline 或 trace drift 超过 threshold 时告警。
\end{knowledgebox}

\subsection{Budget SLO 与警报}
在部署 inference-time compute pipeline 时，常用 SLO 表格列出 token budget、search depth 与 iteration depth 的上限，并配套 “当 average token per query > budget × 1.1” 或 “search branching > 30%” 的警报。Chen 建议在 SLO 里加入“double-check”阶段，用 stepwise verifier 的两个 score 确保 high-stakes query 不会因为 sampling drift 而失灵。

\begin{table}[H]
\centering
\caption{Inference-time SLO 示例}
\begin{tabular}{p{0.32\textwidth}p{0.32\textwidth}p{0.32\textwidth}}
\toprule
Metric & Target & Alert condition \\
\midrule
Token budget per query & $\leq 400$ tokens & 5-minute moving average > 440 tokens \\
Sampling diversity & 80\% unique candidates & 3 consecutive minutes diversity < 60\% \\
Iteration depth & $\leq 2$ feedback rounds & iteration count > 3 for top 1\% queries \\
Verifier score margin & $\geq 0.15$ & margin drop > 0.05 for 100 queries \\
\bottomrule
\end{tabular}
\end{table}

\begin{warningbox}{警报调优建议}
不要让每一次 trace drift 都弹警报，否则会造成 alarm fatigue。Chen 的经验是把误报阈值设得稍高，真正有问题时再调度人工排查，并辅以 “replay” 功能让团队复盘 sampling log。
\end{warningbox}

\subsection{本章小结}
部署 inference-time pipeline 时必须把 observability、SLO、警报与 log replay 联动，才能在 production 层保持 accuracy 和 reliability，同时把 budget 运行在可控范围之内。

